\section{Coboundary expansion}
\label{sec:expansion}

The cohomology calculation has removed the obstruction to coboundary
expansion. We now obtain a uniform cosystolic inequality from the KMS
construction, verify bounded degree and growth, and compare the weights
on the two skeleta.

\subsection{The quantitative input and its hypotheses}

We recall the Kac--Moody--Steinberg (KMS) construction in the case needed
here. The affine diagram $\widetilde A_n$ has type set $S$, and each
proper subdiagram is a disjoint union of paths of type $A$. Its
\emph{spherical local groups} are the corresponding unitriangular groups
$U_J$, $J\subsetneq S$, with their standard inclusions.
Their \emph{universal completion} is the group obtained by taking the
free product of these groups and identifying each $U_J$ with its image
in $U_K$ whenever $J\subseteq K$. This is the KMS group of type
$\widetilde A_n$ \cite[Definition~4.6]{OVB}.

A finite quotient of the completion defines a \emph{KMS complex} if
its map is injective on every spherical local group and preserves their
pairwise intersections; the complex is formed from the images of the
maximal local groups by the common-intersection coset construction
\cite[Definitions~4.8 and~4.11]{OVB}. A diagram on $n+1$ types is
\emph{$n$-classical} when every proper subdiagram has irreducible
components of classical type $A_r$, $B_r$, $C_r$ or $D_r$
\cite[Definition~4.2]{OVB}. Thus $\widetilde A_n$ is $n$-classical.

We use the following part of the theorem of Oppenheim--Valentiner-Branth
\cite[Theorem~1.2]{OVB}. For fixed $n\ge3$ and sufficiently large $q$,
the $(n-1)$-skeleta $Y$ of finite $n$-dimensional, $n$-classical KMS
complexes over $\mathbb F_q$ satisfy, for some $\eta>0$ independent of
the finite quotient,
\begin{equation}\label{eq:cosystolic}
 \norm{\delta_i f}_{Y,i+1}
 \ge\eta\,\dist_{Y,i}(f,Z^i(Y))
 \qquad(0\le i<n-1).
\end{equation}
The weights are exactly \eqref{eq:weights}. We only need this inequality,
not the additional lower bound on nontrivial cosystoles.

The polynomial local groups above also appear in
\cite[Example~4.13]{OVB}, where the finite quotients use irreducible
polynomials. We verify the recognition conditions for our reductions
modulo $t^m$ directly. Our matrix model supplies
a homomorphism from that completion to $U$, and it is onto by
\cref{cor:generation}. The local groups are realized faithfully by
\cref{lem:patterns}; their intersections are the groups belonging to
the intersecting subsets.

The composite map to $G_m$ has both properties by \cref{lem:patterns}.
Indeed the equality of intersections already holds in the universal
completion: a member of two local groups has, in the matrix model,
image in $U_{J\cap K}$, and injectivity on either local group identifies
it with the corresponding member of that common subgroup.
Therefore $T_m$ is an $n$-dimensional, $n$-classical KMS complex.
This also explains why reduction modulo $t^m$ is allowed: the definition
requires a finite group quotient with the two stated local properties,
and imposes no field condition on the quotient ring.

Choose $q$ once and for all to satisfy $q\ge5$, $q\ge2^{n-1}$, and the
largeness requirement for \eqref{eq:cosystolic}. There are arbitrarily
large odd primes. Thus
\[
 Y_m=\skel_{n-1}T_m
\]
satisfies \eqref{eq:cosystolic} with a common constant $\eta>0$.
It remains to pass from these $(n-1)$-skeleta to the $d$-skeleta $X_m$
without losing a uniform expansion constant.

\subsection{Degree bounds and growth}

Recall $M=q^{\binom N2}$ from \eqref{eq:M}. The common stabilizer of all
base vertices is $U_{\varnothing}=\{1\}$, so the chambers of $T_m$ are
in bijection with $G_m$. A face of type $A$ lies in exactly
$|U_{S\setminus A}|$ chambers. In particular, every vertex lies in exactly
$M$ chambers, and every panel (a face of dimension $n-1$) lies in exactly
$q$ chambers. Each chamber has $N$ vertices, so every vertex of $X_m$
lies in at most
\begin{equation}\label{eq:D}
 D=M\binom{N-1}{d}
\end{equation}
faces of dimension $d$. The complexes $T_m$ and $X_m$ are pure, and
$X_m$ is connected by \cref{prop:homology-comparisons}.

The subgroup $\{e_{12}(f):f\in R\}$ lies in $U$ by definition, and its
image in $G_m$ has order $q^m$. There are $|G_m|/M$ vertices of each
of the $N$ types. Thus
\begin{equation}\label{eq:growth}
 |X_m(0)|=\frac{N|G_m|}{M}\ge\frac{Nq^m}{M}\longrightarrow\infty.
\end{equation}

\subsection{Weights on skeleta}

The expansion theorem uses the intrinsic weights of $Y_m$, while
\cref{thm:main} uses those of $X_m$. The following elementary observation
keeps track of this change of normalization.
Let $T$ be a finite pure $n$-complex, put $X=\skel_s T$, and write
$\mu_i=w_i^T$ and $\nu_i=w_i^X$. For $\tau\in T(s)$ let $c(\tau)$
be its number of containing chambers. If $0<a\le c(\tau)\le b$ for all
$\tau\in T(s)$, then, with
\[
 \Lambda=\frac{\binom{n+1}{s+1}|T(n)|}{|T(s)|},
\]
we have
\begin{equation}\label{eq:skeleton-weights}
 \frac{\Lambda}{b}\mu_i(\sigma)\le\nu_i(\sigma)
 \le\frac{\Lambda}{a}\mu_i(\sigma)
 \qquad(0\le i\le s).
\end{equation}
For $i=s$, this follows from $\nu_s(\tau)=\Lambda\mu_s(\tau)/c(\tau)$.
For $i<s$, use the averaging identities
\[
 \mu_i(\sigma)=\binom{s+1}{i+1}^{-1}
       \sum_{\substack{\tau\in T(s)\\\sigma\subseteq\tau}}\mu_s(\tau),
 \qquad
 \nu_i(\sigma)=\binom{s+1}{i+1}^{-1}
       \sum_{\substack{\tau\in T(s)\\\sigma\subseteq\tau}}\nu_s(\tau).
\]
The first is a double count of chamber--face incidences, and the second
is the definition of skeleton weights. Summing the bounds in degree
$s$ proves \eqref{eq:skeleton-weights}. The same inequalities hold for
norms and distances to any common subspace.

Apply this first with $T=T_m$ and $s=n-1$. Every panel lies in $q$
chambers, so $a=b=\Lambda=q$ and the weights of $Y_m$ are exactly
the weights of $T_m$ in their common degrees. Next take $s=d$.
Every $d$-face lies in between $1$ and $M$ chambers, so
\eqref{eq:skeleton-weights} gives
\begin{equation}\label{eq:norm-comparison}
 \frac{\Lambda_m}{M}\norm{f}_{Y_m,i}
 \le\norm{f}_{X_m,i}
 \le\Lambda_m\norm{f}_{Y_m,i}\qquad(0\le i\le d).
\end{equation}
Here $\Lambda_m=\binom N{d+1}|T_m(n)|/|T_m(d)|$ is common to all degrees
and cancels in the expansion inequality below.

\begin{proof}[Proof of \cref{thm:main}]
Take the complexes in \eqref{eq:construction}. We have proved finiteness,
purity, connectedness, the degree bound \eqref{eq:D}, and the growth
\eqref{eq:growth}. For $i<d$, the cochain spaces and coboundary maps
needed to compute $Z^i$ and $B^i$ are identical for $X_m$, $Y_m$, and
$T_m$. By \cref{prop:quotient-homology}, and by connectedness when $i=0$,
\[
 Z^i(Y_m)=B^i(Y_m)=B^i(X_m)\qquad(0\le i<d).
\]
Combining \eqref{eq:cosystolic} and \eqref{eq:norm-comparison} gives
\begin{align*}
 \norm{\delta_i f}_{X_m,i+1}
 &\ge\frac{\Lambda_m}{M}\norm{\delta_i f}_{Y_m,i+1}\\
 &\ge\frac{\Lambda_m\eta}{M}\dist_{Y_m,i}(f,Z^i(Y_m))\\
 &\ge\frac{\eta}{M}\dist_{X_m,i}(f,B^i(X_m)).
\end{align*}
Thus \eqref{eq:goal} holds with $\eps=\eta/M>0$, independent of $m$
and $i$. Reindexing $m=2,3,\ldots$ gives the required sequence.
\end{proof}
